% !TEX root = ../../mixed-sklyanin-cubics.tex
% Current result: 2026-11-22 / v61.
\subsection{Clifford 中心平方による次数零混合類の実積分}
\label{ifc:subsec:clifford-family}

命題\ref{ifh:prop:central-mixed-obstruction}の障害核に現れる
$\operatorname{HH}^2(S_i)_{-2}\otimes Z(S_j)_2$ のうち、
一つの ordered pair に対応する九次元部分は、全係数上の二次代数族へ積分できる。
その構成には、定数型 PBW 因子の普遍係数を、他因子の中心平方へ写す。
中心三次も同じ底変更で持ち上がり、その商は係数環上自由となる。
この族では境界代数自体の次数零変形を行う。
定理\ref{ifp:thm:factorization}で固定した最高 active 次数の tensor 境界も変化している。

\subsubsection{九変数族と中心三次商}

$S_1=S(A_1,B_1,C_1)$ を任意の滑らかな標準 Sklyanin 因子とし、
生成元を $\mathbf x=(x_1,x_2,x_3)$、三つの巡回関係を $r_{1,a}$ と書く。
第二因子を
\[
 S_2=S(1,1,\tau),\qquad
 \tau\ne0,\quad \tau^3\ne1,-8
\]
と固定し、その生成元を $\mathbf y=(y_1,y_2,y_3)$ とする。
このとき
\[
 \{y_2,y_3\}+\tau y_1^2=0,\qquad\text{巡回式},\qquad
 Z(S_2)_2=\langle y_1^2,y_2^2,y_3^2\rangle
\]
である。各因子の指定中心正則三次を $\Omega_{i,0}$ とし、
$a_1,a_2\in\C^\times$ を固定する。
独立な九変数 $\Lambda_{ab}$ の係数環と中心次数二の元を
\begin{equation}\label{ifc:eq:parameters}
 R=\C[\Lambda_{ab}:1\le a,b\le3],\qquad
 z_a=\sum_{b=1}^3\Lambda_{ab}y_b^2\in Z(R\otimes_\C S_2)_2
\end{equation}
と置く。$\Lambda_{ab}$ の内部次数は零である。

\begin{theorem}[中心平方を代入した二次族と自由中心商]
\label{ifc:thm:family}
次の次数付き $R$-代数を考える：
\begin{equation}\label{ifc:eq:presentation}
 \mathcal B_\Lambda
 =R\langle\mathbf x,\mathbf y\rangle
   \big/
    \bigl(r_{1,a}-z_a,\ r_{2,b},\
                   [x_a,y_b]:1\le a,b\le3\bigr).
\end{equation}
全生成元の内部次数を一とする。
この族は $R$ 上自由で、$\Lambda=0$ のファイバーは $S_1\otimes_\C S_2$ である。
各 $S_i$ の固定した斉次語基底の積が $\mathcal B_\Lambda$ の $R$-基底となる。

固定した一次式 $\ell_a(\mathbf x)$ を用いて、中心三次
\begin{equation}\label{ifc:eq:center-lift}
 \Omega_{1,\Lambda}
   =\Omega_{1,0}+\sum_{a,b}\Lambda_{ab}y_b^2\ell_a(\mathbf x)
       \in Z(\mathcal B_\Lambda)_3
\end{equation}
を選べる。ここで $\Omega_{1,0}$ には固定した三次の語代表を用いる。
また $\Omega_{2,0}$ も $\mathcal B_\Lambda$ の中心元である。
従って
\[
 w_\Lambda=a_1\Omega_{1,\Lambda}+a_2\Omega_{2,0},
       \qquad \mathcal A_\Lambda=\mathcal B_\Lambda/(w_\Lambda)
\]
が定まり、$w_\Lambda$ は左右の非零因子である。
$\mathcal A_\Lambda$ も各内部次数で有限自由な $R$-加群であり、Hilbert 級数は
\begin{equation}\label{ifc:eq:hilbert}
 H_{\mathcal B_\Lambda}(t)=\frac1{(1-t)^6},\qquad
 H_{\mathcal A_\Lambda}(t)=\frac{1-t^3}{(1-t)^6}.
\end{equation}
これらの基底・中心性・正則性は、任意の可換 $\C$-代数への係数変更と両立する。
\end{theorem}

\begin{proof}
第一因子の普遍定数型 PBW 代数を
\[
 C_\alpha=\C[\alpha_1,\alpha_2,\alpha_3],\qquad
 \mathcal U_1=C_\alpha\langle\mathbf x\rangle/
                         (r_{1,a}-\alpha_a:1\le a\le3)
\]
とする。$\deg\alpha_a=2$ とし、$\alpha_a$ は中心とする。
一因子の普遍 PBW 基底と中心持上げは、
非二捩れの場合には命題\ref{eco:prop:universal-lift}と
系\ref{eco:cor:non-clifford-lift}、二捩れの場合には
命題\ref{eco:prop:clifford-center-freedom}による。
前者では線形補正を零にして Rees 変数を一へ特殊化すればよい。
従って $\mathcal U_1$ は $C_\alpha$ 上自由で、
$S_1$ の固定した斉次語基底が基底を与える。
また、係数重み二の $\alpha$ と生成元次数一に関する斉次性から、
その中心持上げは
\[
              \Omega_1(\alpha)
                  =\Omega_{1,0}+\sum_a\alpha_a\ell_a(\mathbf x)
\]
と書ける。二捩れの第一因子についても、普遍 Clifford 体積元の同じ重みを用いる。

式\eqref{ifc:eq:parameters}は中心への次数付き準同型
$C_\alpha\to Z(R\otimes S_2)$、$\alpha_a\mapsto z_a$ を与える。
そこで
\begin{equation}\label{ifc:eq:central-base-change}
           \mathcal U_1\otimes_{C_\alpha}(R\otimes_\C S_2)
\end{equation}
を取る。第二因子は非可換だが、$C_\alpha$ は両因子の中心へ写るので
$(p\otimes s)(q\otimes t)=pq\otimes st$ は balanced tensor 積上で定義される。
この代数の普遍性から、式\eqref{ifc:eq:presentation}と同型である。
$\mathcal U_1$ の $C_\alpha$-自由基底を tensor すると主張した基底を得る。
特に $R\otimes S_2$ の $C_\alpha$ 上の平坦性を仮定する必要はない。

$\Omega_1(\alpha)\otimes1$ は第一因子と可換であり、
tensor 積の第二因子とも可換である。その像が
式\eqref{ifc:eq:center-lift}である。
$1\otimes\Omega_{2,0}$ も同様に中心的である。
従って、点ごとに中心元を選び直すことなく、全係数上の中心恒等式を得る。

正則性には、$\mathbf x$ のみを次数一、
$R\otimes S_2$ を次数零とする非負濾過 $F$ を使う。
第一因子の PBW 基底から
\begin{equation}\label{ifc:eq:filtered-algebra}
          \operatorname{gr}_F\mathcal B_\Lambda
                 \simeq S_1\otimes_\C S_2\otimes_\C R
\end{equation}
であり、$w_\Lambda$ の最高 $F$-次数部分は
$a_1\Omega_{1,0}\otimes1$ である。
$\Omega_{1,0}$ は $S_1$ 上の中心正則元であるため、
この最高次部分は式\eqref{ifc:eq:filtered-algebra}上で左右正則である。
有限な最高次数を持つ元の積の初項を比較すれば、
$w_\Lambda$ 自身も左右正則と分かる。

さらに、主イデアルの濾過も strict である。
実際、非零元 $b$ に対し
$\deg_F(w_\Lambda b)=3+\deg_F b$ が成り立ち、
中心性によって二側イデアルの全元は $w_\Lambda b$ の形に書ける。
従って
\begin{equation}\label{ifc:eq:filtered-quotient}
 \operatorname{gr}_F\mathcal A_\Lambda
        \simeq (S_1/(\Omega_{1,0}))\otimes_\C S_2\otimes_\C R.
\end{equation}
$S_1/(\Omega_{1,0})$ の斉次語基底を持ち上げたものと、
$S_2$ の斉次語基底との積を取る。
式\eqref{ifc:eq:filtered-quotient}と最高 $F$-次数に関する帰納法により、
それらは $\mathcal A_\Lambda$ の $R$-自由基底となる。
全ての関係と $w_\Lambda$ は通常の内部次数について斉次なので、
この基底は内部次数別にも斉次である。
各 Sklyanin 因子の Hilbert 級数と $\deg w_\Lambda=3$ から
式\eqref{ifc:eq:hilbert}を得る。
全ての基底表示と初項の単射性は $\C$ 上の自由空間を用いているため、
非被約な係数環への変更にも同じ証明が適用できる。
\end{proof}

\subsubsection{接類と一次障害核への対応}

\begin{proposition}[九つの混合方向の独立性]
\label{ifc:prop:tangent}
$T=S_1\otimes S_2$ とし、二因子次数を
$\deg\mathbf x=(1,0)$、$\deg\mathbf y=(0,1)$ とする。
族 $\mathcal B_\Lambda$ の原点における接写像は同型
\begin{equation}\label{ifc:eq:tangent}
 \operatorname{Mat}_3(\C)
   \xrightarrow{\ \sim\ }
 \operatorname{HH}^2(T)_{(-2,2)}
   \simeq \operatorname{HH}^2(S_1)_{-2}\otimes Z(S_2)_2
       \subset\operatorname{HH}^2(T)_0
\end{equation}
を与える。定数補正 $r_{1,a}-\alpha_a$ に対応する類を $\eta_a$ とすれば、
行列 $M$ の像は $\sum_{a,b}M_{ab}\eta_a\otimes y_b^2$ である。
従って九方向は全て非零かつ独立であり、
命題\ref{ifh:prop:central-mixed-obstruction}の
$\ker\operatorname{ob}_{w_0}$ に属する。
本族はこの一つの ordered pair の九方向を同時に全次数へ積分する。
\end{proposition}

\begin{proof}
式\eqref{ifh:eq:kunneth}の Koszul 双加群複体を二因子次数で分解する。
$R_1$ を第一因子の関係空間とすると、次数の非負性から
\[
 C^1_{(-2,2)}=0,\qquad
 C^2_{(-2,2)}=\operatorname{Hom}_\C(R_1,(S_2)_2)
\]
である。後者の写像を $\psi$ とする。
第一因子だけの三次 overlap では、cocycle 条件は
$\sum_a[x_a,\psi(r_{1,a})]=0$ であり、異なる因子の可換性から自動的に成立する。
二次の第一因子と一次の第二因子の shuffle overlap では、
条件は
\[
                [y,\psi(r)]=0
             \qquad(y\in(S_2)_1,\ r\in R_1)
\]
である。他の三次成分は第一因子の出力次数が負となり寄与しない。
従って cocycle 空間は正確に
$\operatorname{Hom}_\C(R_1,Z(S_2)_2)$ となる。
$C^1_{(-2,2)}=0$ により境界はなく、
三つの関係双対と三つの独立な中心平方から次元は九である。
これは純因子だけでなく、全 Hochschild cochain に対する計算である。

独立パラメータ族を双対数へ $\Lambda_{ab}=\epsilon M_{ab}$、
$\epsilon^2=0$ と制限する。
定理\ref{ifc:thm:family}の自由基底で一次の乗法を取り出すと、
Koszul 二 cochain の $r_{1,a}$ 上の値は
$\sum_bM_{ab}y_b^2$ であり、他の関係上では零である。
これが式\eqref{ifc:eq:tangent}の同型を与える。
ここで接点は九変数 $\Lambda$ の原点であり、$M$ はその接ベクトルである。
最後に、全係数上で存在する $w_\Lambda$ を同じ双対数へ制限すれば、
$w_0=a_1\Omega_{1,0}+a_2\Omega_{2,0}$ の中心持上げを得る。
従って一次中心障害は零であり、実際にはこの族で全次数の持上げが構成されている。
\end{proof}

\begin{corollary}[自由な中心変数の追加]
\label{ifc:cor:central-variables}
任意個の中心次数一変数 $\mathbf u=(u_1,\ldots,u_m)$ を自由に追加し、
$g(\mathbf u)$ を任意の斉次三次式とする。
このとき
\[
 \mathcal B_\Lambda[\mathbf u],\qquad
 \mathcal A_{\Lambda,g}
   =\mathcal B_\Lambda[\mathbf u]/
       (w_\Lambda+g(\mathbf u))
\]
についても、表示した中心三次は左右正則であり、
$\mathcal A_{\Lambda,g}$ は $R[\mathbf u]$ 上自由である。
各複素係数点での内部 Hilbert 級数は
\[
              H_{\mathcal A_{\Lambda,g}}(t)
                       =\frac{1-t^3}{(1-t)^{6+m}}
\]
である。$g$ の係数も独立な変数として加えてよい。
\end{corollary}

\begin{proof}
前証明の $F$ に関して $\mathbf u$ の次数を零とする。
中心三次の最高 $F$-次数部分は引き続き $a_1\Omega_{1,0}$ である。
式\eqref{ifc:eq:filtered-quotient}の右辺に
$\C[\mathbf u]$ と $g$ の係数環を tensor し、同じ基底の議論を適用すればよい。
\end{proof}

\begin{remark}[この積分族の範囲]
\label{ifc:rem:scope}
本小節で得たものは、固定 Clifford 因子の中心平方を使う、
一方向の九変数混合族とその中心正則三次商である。
式\eqref{ifh:eq:obstruction-kernel}の他の ordered pair、
逆向きの混合類、純因子の次数零変形と同時に積分できるとは結論していない。
また、この $F$ は第一因子だけの次数濾過であり、
固定 tensor 境界の中心一次・二次補正を扱う前小節の active 濾過とは異なる。
PBW 基底、係数環上の自由性、中心持上げとその正則性がここでの結論である。
周囲代数の AS-regular 性・CY 性、射影中心商や残余圏の smoothness、
および可換三次への Ku 比較核は本小節では確認していない。
中心変数を加えない Hilbert 型は四次元三次に対応するため、
主稿の $n\ne4$ という仮定を持つ開性定理
（定理\paperref{main}{grl:thm:open-locus}）をそのまま用いることもできない。
一つの中心変数と非零三乗項を加えた五次元型では、この次元上の制約はなくなるが、
同定理に必要な実 smooth proper 圏族などの入力は別途確認しなければならない。
\end{remark}
